\chapter{Catálogo de Estrategias de Mitigación}
\label{anexo:catalogo-estrategias-mitigacion}

Este anexo presenta un catálogo de estrategias de mitigación para diferentes categorías de \textit{community smells}. Cada estrategia incluye una descripción, pasos recomendados y consideraciones clave para su implementación efectiva.

\section{Estrategias de mitigación de community smells}

Detectar un \textit{community smell} es solo el primer paso. La detección sin acción genera frustración: el equipo sabe que hay un problema pero no ve mejora. Esta sección proporciona un catálogo de estrategias de mitigación basadas en la taxonomía de \textit{community smells} de Caballero et al.\ \cite{caballero2022communitysmells}, organizadas por categoría de smell. Herramientas como csDetector \cite{almarimi2021csDetector} permiten detección automática de algunos de estos patrones a partir de metadatos de repositorios.

El objetivo no es eliminar todos los smells ---algunos son síntomas de tensiones organizacionales que el equipo no controla--- sino reducir su impacto en la capacidad del equipo de entregar valor.

\subsection{Catálogo de estrategias por categoría}

Para cada categoría de \textit{community smell}, se describen problemas comunes, ejemplos representativos y tipos de intervención recomendados (ver Tabla~\ref{tab:categorias-smells}).

\begin{table}[!htbp]
    \centering
    \small
    \begin{tabularx}{\textwidth}{|l|X|X|X|}
        \hline
        \textbf{Categoría} & \textbf{Problema}                                                                                                                      & \textbf{Smells representativos} & \textbf{Tipo de intervención} \\
        \hline
        Comunicación
                           & Información no fluye entre miembros o hacia afuera
                           & \textit{Radio Silence}, \textit{Information Hoarding}, \textit{Broadcast Storm}, \textit{Missing Links}, \textit{Synchronous Overload}
                           & Rediseño de canales, protocolos y automatización                                                                                                                                                         \\
        \hline
        Conocimiento
                           & Saber técnico concentrado o no compartido
                           & \textit{Truck Factor} = 1, \textit{Knowledge Islands}, \textit{Outdated Docs}, \textit{Tribal Knowledge}, \textit{Expert Bottleneck}
                           & \textit{Pair programming}, documentación y rotación                                                                                                                                                      \\
        \hline
        Relacional
                           & Conflictos, exclusión y dinámicas tóxicas
                           & \textit{Lone Wolf}, \textit{Prima Donna}, \textit{Clique Formation}, \textit{Unresolved Conflict}, \textit{Toxic Leadership}
                           & Mediación, facilitación y feedback directo                                                                                                                                                               \\
        \hline
        Organizacional
                           & Estructura o procesos disfuncionales
                           & \textit{Organizational Silo}, \textit{Meeting Hell}, \textit{Unclear Ownership}, \textit{Process Overhead}, \textit{Unstable Team}
                           & Escalamiento, negociación y reestructuración                                                                                                                                                             \\
        \hline
    \end{tabularx}
    \caption{Categorías de \textit{community smells}, ejemplos e intervenciones}
    \label{tab:categorias-smells}
\end{table}
\FloatBarrier

\subsubsection{Estrategias para \glsentrytext{community-smells} de comunicación}

Estas estrategias abordan problemas de flujo de información dentro del equipo y con \textit{stakeholders} externos.

\begin{table}[!htbp]
    \centering
    \footnotesize
    \begin{tabularx}{\textwidth}{|l|X|X|c|}
        \hline
        \textbf{Smell} & \textbf{Síntomas}                                                       & \textbf{Estrategia de mitigación} & \textbf{SSP} \\
        \hline
        \textit{Radio Silence}
                       & Sorpresas frecuentes; “nadie me dijo”; la información llega tarde
                       & 1) Identificar nodos centrales mediante SNA.
        2) Crear un canal de broadcast para actualizaciones críticas.
        3) Establecer \textit{office hours} donde cualquiera pueda preguntar.
        4) Rotar la responsabilidad de comunicar en la Daily.
                       & 3--5                                                                                                                       \\
        \hline
        \textit{Information Hoarding}
                       & Personas guardan información “por si acaso”; duplicación de esfuerzo
                       & 1) Workshop sobre beneficios de compartir.
        2) Crear una wiki o Confluence con expectativa explícita de documentar decisiones.
        3) Reconocer públicamente cuando alguien comparte información útil.
                       & 3--5                                                                                                                       \\
        \hline
        \textit{Broadcast Storm}
                       & Canales saturados; nadie lee mensajes; información importante se pierde
                       & 1) Auditar y consolidar canales existentes.
        2) Establecer convenciones (por ejemplo, @here solo para urgencias).
        3) Crear resúmenes diarios o semanales.
        4) Eliminar canales redundantes.
                       & 2--3                                                                                                                       \\
        \hline
        \textit{Missing Links}
                       & Subgrupos no se comunican; integraciones tardías generan conflictos
                       & 1) Mapear la comunicación actual mediante SNA.
        2) Crear rituales de sincronización entre subgrupos.
        3) Asignar enlaces rotativos entre equipos.
        4) Realizar retrospectivas conjuntas mensuales.
                       & 5--8                                                                                                                       \\
        \hline
        \textit{Synchronous Overload}
                       & Calendarios saturados; trabajo fragmentado; fatiga de videollamadas
                       & 1) Auditar reuniones y eliminar las innecesarias.
        2) Establecer \textit{no-meeting days}.
        3) Convertir reuniones de estado a formato asíncrono.
        4) Limitar la duración máxima de reuniones (25/50 min).
                       & 3--5                                                                                                                       \\
        \hline
    \end{tabularx}
    \caption{Estrategias para smells de comunicación}
    \label{tab:estrategias-comunicacion}
\end{table}
\FloatBarrier

\subsubsection{Estrategias para smells de conocimiento}

Estas estrategias se enfocan en mejorar la gestión y transferencia del conocimiento dentro del equipo.

\begin{table}[!htbp]
    \centering
    \footnotesize
    \begin{tabularx}{\textwidth}{|l|X|X|c|}
        \hline
        \textbf{Smell} & \textbf{Síntomas}                                                              & \textbf{Estrategia de mitigación} & \textbf{SSP} \\
        \hline
        \textit{Truck Factor} = 1
                       & “Solo Juan sabe cómo funciona X”; bloqueos cuando Juan no está
                       & 1) Identificar componentes críticos con un solo experto.
        2) \textit{Pair programming} forzado: Juan + otro dev en cada tarea del componente.
        3) Juan documenta la arquitectura básica (2–3 páginas).
        4) Otro dev realiza un cambio pequeño de forma autónoma; Juan revisa.
                       & 5--8                                                                                                                              \\
        \hline
        \textit{Knowledge Islands}
                       & Subgrupos especializados no comparten conocimiento; la integración es dolorosa
                       & 1) \textit{Tech talks} semanales rotativos (30 min).
        2) \textit{Code reviews} cruzados entre subgrupos.
        3) Rotación temporal: un Sprint en otro subgrupo.
        4) Documentación explícita de interfaces entre componentes.
                       & 5--8                                                                                                                              \\
        \hline
        \textit{Outdated Documentation}
                       & La documentación existe pero es incorrecta; nadie confía en ella
                       & 1) Auditar documentación: marcar obsoleta o eliminar.
        2) Establecer \textit{docs as code}: la documentación vive junto al código.
        3) Incluir actualización de documentación en la \textit{Definition of Done}.
        4) Asignar un responsable de cada documentación crítica.
                       & 3--5                                                                                                                              \\
        \hline
        \textit{Tribal Knowledge}
                       & Conocimiento principalmente oral; el \textit{onboarding} toma meses
                       & 1) Grabar sesiones de \textit{onboarding} para reutilizarlas.
        2) Crear un \textit{playbook} de operaciones comunes.
        3) \textit{Pair programming} intensivo durante las primeras dos semanas.
        4) Implementar un sistema de compañeros (\textit{buddy system}) para nuevos miembros.
                       & 5--8                                                                                                                              \\
        \hline
        \textit{Expert Bottleneck}
                       & Una sola persona revisa todo; los PRs esperan días
                       & 1) El experto entrena a 2–3 personas como reviewers secundarios.
        2) Establecer un SLA (Service Level Agreement): PRs simples no requieren aprobación del experto.
        3) Crear checklists de revisión reutilizables.
        4) Delegar gradualmente categorías de PRs.
                       & 5--8                                                                                                                              \\
        \hline
    \end{tabularx}
    \caption{Estrategias para smells de conocimiento}
    \label{tab:estrategias-conocimiento}
\end{table}
\FloatBarrier

\subsubsection{Estrategias para smells relacionales}

Estas estrategias se enfocan en mejorar las interacciones y relaciones dentro del equipo.

\begin{table}[!htbp]
    \centering
    \footnotesize
    \begin{tabularx}{\textwidth}{|l|X|X|c|}
        \hline
        \textbf{Smell} & \textbf{Síntomas}                                                                     & \textbf{Estrategia de mitigación} & \textbf{SSP} \\
        \hline
        Lone Wolf
                       & Trabaja solo; no pide ni ofrece ayuda; PRs sin contexto
                       & 1) 1:1 para entender la motivación (preferencia o desconfianza).
        2) Pair programming gradual (1 hora al día).
        3) Asignar tareas que requieran colaboración.
        4) Reconocer públicamente cuando colabora.
                       & 3--5                                                                                                                                     \\
        \hline
        Prima Donna
                       & Cree que sus ideas son superiores; no acepta feedback; bloquea el consenso
                       & 1) 1:1 para dar feedback específico sobre el impacto.
        2) Establecer la norma de decisiones por consenso, no por imposición.
        3) Documentar criterios objetivos para decisiones técnicas.
        4) Si persiste, escalar al Scrum Master.
                       & 5--8                                                                                                                                     \\
        \hline
        Clique Formation
                       & Subgrupo cerrado; exclusión en almuerzos o canales; la información no sale del clique
                       & 1) Workshop de cohesión con todo el equipo.
        2) Rotar parejas de trabajo para romper cliques.
        3) Crear rituales inclusivos (coffee chat aleatorio).
        4) 1:1 con líderes informales del clique.
                       & 5--8                                                                                                                                     \\
        \hline
        Unresolved Conflict
                       & Tensión visible; evitan trabajar juntos; comentarios pasivo-agresivos
                       & 1) 1:1 con cada parte para entender su perspectiva.
        2) Sesión de mediación facilitada.
        3) Documentar acuerdos de comportamiento.
        4) Seguimiento en 2--4 semanas.
                       & 5--8                                                                                                                                     \\
        \hline
        Toxic Leadership
                       & Micromanagement; culpabilización pública; generación de miedo
                       & 1) Documentar incidentes específicos con fechas.
        2) Escalar a CTO o VP con evidencia.
        3) Proteger al equipo mientras se resuelve la situación.
        4) En casos severos, separar al líder del equipo.
                       & 8--13                                                                                                                                    \\
        \hline
    \end{tabularx}
    \caption{Estrategias para smells relacionales}
    \label{tab:estrategias-relacionales}
\end{table}
\FloatBarrier

\subsubsection{Estrategias para smells organizacionales}

Estas estrategias abordan problemas estructurales y de proceso que afectan al equipo.

\begin{table}[!htbp]
    \centering
    \footnotesize
    \begin{tabularx}{\textwidth}{|l|X|X|c|}
        \hline
        \textbf{Smell} & \textbf{Síntomas}                                                      & \textbf{Estrategia de mitigación} & \textbf{SSP} \\
        \hline
        \textit{Organizational Silo}
                       & Equipos no se comunican; duplicación de trabajo; integraciones tardías
                       & 1) Identificar puntos de contacto necesarios.
        2) Establecer embajadores entre equipos.
        3) Reuniones de sincronización bisemanales.
        4) Escalar a liderazgo si la estructura impide la colaboración.
                       & 5--8                                                                                                                      \\
        \hline
        \textit{Meeting Hell}
                       & Calendarios saturados por demandas externas al equipo
                       & 1) Auditar reuniones: quién las convoca y si son necesarias.
        2) Negociar con \textit{stakeholders} un representante único.
        3) Establecer \textit{focus time} protegido.
        4) Escalar al Scrum Master para proteger al equipo.
                       & 3--5                                                                                                                      \\
        \hline
        \textit{Unclear Ownership}
                       & Preguntas frecuentes sobre responsables; decisiones postergadas
                       & 1) Mapear componentes y asignar responsables explícitos.
        2) Documentar una matriz RACI para decisiones comunes.
        3) Publicar el \textit{ownership} en una wiki accesible.
        4) Revisar y actualizar cada tres meses.
                       & 3--5                                                                                                                      \\
        \hline
        \textit{Process Overhead}
                       & Aprobaciones múltiples y burocracia que no agrega valor
                       & 1) Documentar tiempo perdido en procesos.
        2) Proponer simplificación basada en evidencia.
        3) Ejecutar un piloto con exención de un proceso durante dos Sprints.
        4) Medir impacto y escalar resultados.
                       & 5--8                                                                                                                      \\
        \hline
        \textit{Unstable Team}
                       & Miembros prestados y rotación frecuente
                       & 1) Documentar el impacto en velocidad y calidad.
        2) Escalar a liderazgo con datos.
        3) Negociar un período mínimo de estabilidad.
        4) Si no es posible cambiar, optimizar el onboarding.
                       & 5--8                                                                                                                      \\
        \hline
    \end{tabularx}
    \caption{Estrategias para smells organizacionales}
    \label{tab:estrategias-organizacionales}
\end{table}
\FloatBarrier

\subsubsection{Smells que requieren escalamiento}

Algunos smells no pueden resolverse a nivel de equipo. El Social Guide debe reconocer estos casos y escalar apropiadamente:

\begin{table}[!htbp]
    \centering
    \small
    \begin{tabularx}{\textwidth}{|l|X|X|}
        \hline
        \textbf{Smell} & \textbf{Por qué requiere escalamiento}                                            & \textbf{A quién escalar} \\
        \hline
        \textit{Toxic Leadership}
                       & El líder tiene poder formal sobre el equipo; la confrontación directa es riesgosa
                       & CTO, VP Engineering, HR (con cuidado)                                                                        \\
        \hline
        \textit{Organizational Silo} (estructural)
                       & Equipos separados por decisiones organizacionales
                       & VP Engineering, Director de Producto                                                                         \\
        \hline
        \textit{Unstable Team} (política)
                       & Las decisiones de staffing provienen de niveles superiores
                       & CTO, Resource Manager                                                                                        \\
        \hline
        \textit{Budget Constraints}
                       & Problemas derivados de falta de recursos
                       & CTO, VP con autoridad presupuestaria                                                                         \\
        \hline
        Legal / Problemas de cumplimiento (\textit{Compliance Issues})
                       & Procesos impuestos por regulación o cumplimiento normativo
                       & Legal, Compliance Officer                                                                                    \\
        \hline
    \end{tabularx}
    \caption{Smells que requieren escalamiento}
    \label{tab:smells-escalamiento}
\end{table}
\FloatBarrier

Para estos casos, el rol del Social Guide es:
\begin{enumerate}
    \item Documentar el impacto con datos cuantitativos.
    \item Preparar propuesta de solución (no solo el problema).
    \item Escalar al \textit{stakeholder} apropiado con Scrum Master.
    \item Proteger al equipo mientras se resuelve externamente.
\end{enumerate}

\subsection{Diseño de experimentos de mejora}

Las estrategias del catálogo son puntos de partida que pueden requerir adaptaciones al contexto. El diseño de experimentos permite probar intervenciones sistemáticamente.

\begin{table}[!htbp]
    \centering
    \small
    \begin{tabularx}{\textwidth}{|l|X|}
        \hline
        \textbf{Componente} & \textbf{Descripción}                                                                        \\
        \hline
        Hipótesis
                            & Declaración falsificable que conecta una intervención con un resultado esperado.
        Formato: \textit{Si [intervención], entonces [resultado], porque [mecanismo]}.                                    \\
        \hline
        Métricas
                            & Indicadores cuantitativos definidos con baseline y objetivo (\textit{target}).              \\
        \hline
        Diseño
                            & Definición de qué se hará, cuándo, quiénes participan y recursos requeridos (SSP).          \\
        \hline
        Criterios de evaluación
                            & Éxito: métricas alcanzan el \textit{target} $\rightarrow$ institucionalizar. \newline
        Parcial: mejora $\geq$ 50\% $\rightarrow$ iterar. \newline
        Fracaso: sin mejora $\rightarrow$ pivotar. \newline
        Inconcluso: datos insuficientes $\rightarrow$ extender experimento.                                               \\
        \hline
        Documentación
                            & Registro estructurado de la hipótesis, resultados obtenidos, aprendizajes y decisión final. \\
        \hline
    \end{tabularx}
    \caption{Componentes de un experimento de mejora}
    \label{tab:experimento-mejora}
\end{table}
\FloatBarrier

No toda intervención requiere experimento formal. Se recomienda diseño experimental cuando: el contexto es único, la intervención es costosa ($>$8 SSP), el equipo es escéptico, o el impacto afecta múltiples equipos. Intervenciones simples y reversibles pueden implementarse directamente y ajustarse.